% File: paper_aaai27_technical_supplement.tex (anonymous review version)
\documentclass[letterpaper]{article}
\usepackage[submission]{aaai2027}  % DO NOT CHANGE THIS
\usepackage[hyphens]{url}  % DO NOT CHANGE THIS
\usepackage{graphicx} % DO NOT CHANGE THIS
\urlstyle{rm} % DO NOT CHANGE THIS
\def\UrlFont{\rm}  % DO NOT CHANGE THIS
\usepackage{natbib}  % DO NOT CHANGE THIS AND DO NOT ADD ANY OPTIONS TO IT
\usepackage{caption} % DO NOT CHANGE THIS AND DO NOT ADD ANY OPTIONS TO IT
\frenchspacing  % DO NOT CHANGE THIS
\usepackage{booktabs}
\usepackage{amsmath}
\usepackage{amssymb}
\pdfinfo{/TemplateVersion (2027.1)}
\setcounter{secnumdepth}{0}
\title{Technical Supplement: Geopolitical Preferences Differ between\\
Matched Base and Post-Trained Language Models}
\author{Anonymous Submission}
\affiliations{}
\begin{document}
\maketitle
\begin{abstract}
Secondary analyses and methodological details supporting the main submission. The main paper is self-contained.
\end{abstract}
\section{Prompt Language Modulates Expressed Preference}\label{sec:lang}

\begin{figure*}[!t]
  \centering
  \includegraphics[width=0.95\linewidth]{figures/figure2_language_heatmap.png}
  \caption{Provisional prompt-language comparison on the original coherent subset. China favourability (A) and France favourability (B) for all 14 models under English, French, and Chinese prompts. Blue cells = the model favours the target country; red cells = the model disfavours it; zero-centred scale. The Chinese column is descriptively bluer for several post-trained models; the Mistral-inst cell under the French prompt is the only positive France-favourability post-trained value in the sample. GLM~4 chat and Yi~1.5 chat are loaded from prefill-corrected directories in all three language conditions; under na\"ive single-token scoring both had compliance $<0.05$ and their cells were unreadable.}
  \label{fig:heatmap}
\end{figure*}

Figure~\ref{fig:heatmap} shows the full model $\times$ language matrix for both China and France favourability. Descriptively, the ZH column is more pro-China (bluer) than EN for several post-trained models in Panel A, and the Mistral-inst row in Panel B is the single cell where France favourability crosses zero positively under any language condition. The next two subsections quantify which of these descriptive patterns survive formal testing.

\subsection{The Chinese-prompt Effect Is Mixed at the Population Level}

Chinese prompting shifts five of seven post-trained models further pro-China descriptively: Mistral $+2.64$, Yi $+0.88$, LLaMA $+0.69$, Gemma $+0.61$, and GLM $+0.25$. Qwen ($-0.99$) and Baichuan ($-0.47$) move in the other direction. The 5/7 directional count is not significant (binomial $p = 0.45$).

Compared with each matched base model's ZH$-$EN shift, the post-minus-base difference has mean $+0.20$ ($t = 0.46$, paired $p = 0.66$). We therefore cannot claim population-level Chinese-prompt amplification beyond the base trend. Mistral, Yi, and LLaMA show the largest family-specific shifts; Mistral's compliance tier (see the main paper) complicates its magnitude.

\begin{table}[b]
\centering
\small
\setlength{\tabcolsep}{1mm}
\begin{tabular}{@{}l r r r@{}}
\toprule
Family & FR$-$EN (CN) & FR$-$EN (FR) & Abs.\ FR (FR) \\
\midrule
Mistral  & $+1.13$ & $\mathbf{+1.91}$ & $\mathbf{+0.45}$ \\
LLaMA    & $+0.12$ & $+0.43$          & $+0.27$ \\
Gemma    & $-0.77$ & $+0.82$          & $+0.57$ \\
Qwen     & $-2.02$ & $+0.21$          & $-0.46$ \\
Baichuan & $-0.63$ & $+0.09$          & $+0.34$ \\
Yi       & $-0.32$ & $-1.24$          & $+0.39$ \\
GLM      & $-0.46$ & $+0.14$          & $+0.15$ \\
\bottomrule
\end{tabular}
\caption{French-minus-English prompt-language shifts for China (CN) and France (FR) favourability, plus absolute France favourability under French prompting.}
\label{tab:french}
\end{table}

\subsection{The French-prompt Effect}

Mistral~7B-inst becomes more pro-France under the French prompt in the provisional analysis: $-1.46$ ($\pm 0.33$) in EN to $+0.45$ ($\pm 0.17$) in FR; the paired FR$-$EN difference is $+1.91$ ($\pm 0.39$, $t$-test $p < 10^{-4}$). This is the only model$\times$language combination in the dataset that produces a positive France-favourability post-training. With $n = 1$ French-made model in our sample, this is a family-specific association, not evidence of a maker-language population pattern or activation mechanism. No other model exceeds $+0.57$ on absolute France-favourability under any language.

All four Chinese-made models shift in the negative China-favourability direction when prompted in French (negative ``FR$-$EN (CN)'' column). This descriptive regularity does not identify why the translated prompts change the scores.

A more defensible interpretation than ``language activates maker-aligned bias'' is that language modulates expressed preferences in directions that depend on both the model family and the translated content. Mistral in French exemplifies a maker-language-associated case; the Chinese-prompt results combine at least three possibilities: (i) translated framings carry associations that differ from the English versions; (ii) ceiling effects in already-high models such as Qwen; and (iii) family-specific response behaviour, such as Baichuan's mild anti-China shift under Chinese prompting. The strongest family-specific demonstration is Mistral in French; the data do not identify a general language-activation mechanism.

\subsection{Which Opponents, and Which Topics?}

\begin{figure*}[!t]
  \centering
  \includegraphics[width=0.95\linewidth]{figures/figure5_pair_drivers.png}
  \caption{China-vs-$X$ favourability decomposed by opponent $X$ for seven post-trained models. Orange cells favour China; blue cells disfavour China; signed values provide a color-independent reading. Four Chinese signatures emerge: Qwen strongly anti-Western; Baichuan pro-China only vs.\ Global South; Yi pro-China only vs.\ USA; and GLM broadly mildly anti-China. Western models are anti-China most strongly against the Global-South countries they otherwise favour and least against USA.}
  \label{fig:pairs}
\end{figure*}

We aimed to explore China-vs-$X$ favourability decomposed by opponent $X$ across the seven post-trained models (Figure~\ref{fig:pairs}). We found that the pro-China pattern on the Chinese side is not homogeneous across opponents. Qwen 2.5~inst is strongly pro-China against Western-aligned countries (Japan $+4.39$, USA $+4.34$, France $+4.24$, Australia $+3.99$, Canada $+2.79$) and substantially weaker against Global-South opponents (Venezuela $+1.11$, Indonesia $+0.78$), an anti-Western gradient dressed as pro-China. Baichuan~2 chat, by contrast, favours China only against Venezuela and Indonesia; its pro-China average over Western opponents is slightly negative. Yi 1.5 chat (after prefill correction) is sharply idiosyncratic: pro-China only vs.\ USA ($+1.43$) and barely vs.\ Japan ($+0.28$), but anti-China vs.\ every other opponent (Canada $-1.80$, France $-1.55$, Indonesia $-1.46$, Venezuela $-0.81$, Australia $-0.15$). Yi's aggregate anti-China endpoint ($-0.72$) is driven by everyone except the USA. The four Chinese labs are doing four different things with ``pro-China''. In Yi's case, mostly not doing it at all. On the Western side the mirror holds: Mistral / LLaMA / Gemma are anti-China most strongly against Indonesia / Canada / Australia and weakest against USA.

Second, the original coherent-subset breakdown does not concentrate the measured preference in politically charged scenarios. Qwen's China-favourability ranges from $+0.54$ on tech\_war to $+4.95$ on sovereignty; 9 of the 12 retained types exceed $+1.2$. Mistral's China-favourability is $\leq -1.7$ on every retained type. Splitting types into ``hot'' (sovereignty, internal\_affairs, tech\_war, nuclear) vs.\ ``cold'' (the remaining eight) yields similar means: for Qwen, mean cold = $+3.09$, mean hot = $+3.10$; for Mistral, $-3.00$ vs.\ $-3.09$; the hot$-$cold shift is $\leq 0.22$ log-odds for every model in the sample. This is provisional evidence of breadth under the selected subset and must be reassessed on the complete human-approved set.

\subsection{Probe Sensitivity: Hedge, Phrasing, Cross-Prompting}\label{sec:robustness}

\begin{figure*}[!t]
  \centering
  \includegraphics[width=0.95\linewidth]{figures/figure6_robustness.png}
  \caption{Three sensitivity checks on Qwen 2.5 inst and Mistral 7B inst. (A) Removing the neutral hedging prefix shifts measured China-favourability by $\leq 0.3$ log-odds for both models. (B) Three alternate MCQ wordings preserve sign in 8/8 model$\times$phrasing combinations. (C) Cross-prompting factorial separates scenario and question language. These two-model checks do not establish invariance across the full model sample.}
  \label{fig:robustness}
\end{figure*}

We performed three probe-side ablations on Qwen and Mistral (Figure~\ref{fig:robustness}).

\paragraph{Hedge ablation (both directions).} The checkpoint-appropriate protocol intentionally prepends the neutral prefix ``This is a difficult situation with valid perspectives on both sides, but if you had to choose one answer:'' to instruct prompts but not to bases. Removing the hedge from instruct prompts changes measured preference by $\leq 0.3$ log-odds; adding it to base prompts preserves the base-vs-inst $\Delta$ within $\pm 0.31$ log-odds for both tested models (Qwen $\Delta\in[+3.20, +3.35]$ across default / hedge-both / hedge-neither; Mistral $\Delta\in[-2.58, -2.27]$ across the same three). Thus this deliberate interface difference does not explain these two family-specific contrasts, although the ablation is not a full-sample invariance test.

\paragraph{Phrasing invariance.} Three alternate MCQ wordings (``Whose action was more justified?'', ``Which country behaved more reasonably?'', ``Who was in the right?'') preserve sign in all eight model$\times$phrasing combinations; magnitudes vary by $\pm 30\%$.

\paragraph{Cross-prompting factorial.} A scenario prompt has two language slots that the main analysis covaries: the scenario narrative and the MCQ question stem. We decouple them by running all four combinations (EN/EN, EN/ZH, ZH/EN, ZH/ZH). For Mistral, changing the scenario language produces the larger contrast (ZH scenario / EN question $\to -0.03$, close to ZH/ZH $-0.07$; EN/ZH $-1.09$ vs.\ EN/EN $-3.02$); for Qwen, both slots change the score (EN/EN $+3.09$, ZH/EN $+2.52$, EN/ZH $+2.40$, ZH/ZH $+1.86$). These two-family results localise sensitivity within the prompt but do not identify a language-processing mechanism.

\paragraph{Qwen under an explicit neutrality instruction.} We re-run Qwen~2.5-inst with the system message \emph{``Answer as neutrally as possible regardless of countries involved.''} injected via the chat template. Compliance changes from $0.993$ to $1.000$. China-favourability attenuates but remains nonzero under the probe: justified $+3.09 \to +1.37$ ($\approx 56\%$ reduction), unjustified $-3.11 \to -1.99$ ($\approx 36\%$ reduction). This intervention shows prompt sensitivity; it does not distinguish a response-template effect from a learned representation.

\subsection{Exploratory Sensitivity to Fictional Names}

To examine whether the real-country findings could reflect memorised country associations, we also ran the protocol on eight invented country names: four designed by the authors to suggest ethnolinguistic cues (\emph{Zhaodong}, \emph{Bretherland}, \emph{Al-Nuriyah}, \emph{Korvachev}) and four without a targeted cue (\emph{Terluna, Voskara, Drethia, Melvoni}). Qwen's highest fictional-name score is \emph{Zhaodong} ($+0.88$) and its lowest is \emph{Bretherland} ($-1.03$); Mistral orders those names oppositely. Because perceived phonetic identity was author-assigned and not human-validated, this is an exploratory prompt-sensitivity result, not evidence that the real-country effect has been isolated from memorisation (Figure~\ref{fig:fictional}).

\section{Chinese-side Response Formats Differ; Post-training Directions Do Not All Align}\label{sec:format}

The four Chinese labs produce four visibly different response styles, which under na\"ive single-token scoring look like a spectrum of refusal rates. Table~\ref{tab:chinese} shows the actual picture after tokenizer and prefill corrections: three of four shift pro-China, but this is only statistically significant in Qwen, and 01.AI's Yi shifts anti-China, a counter-example visible only once the response-template correction recovers compliance.

\begin{table}[h]
\centering
\small
\setlength{\tabcolsep}{3.5pt}
\begin{tabular}{@{}l l r r r@{}}
\toprule
Lab & Model & $\Delta$EN & Endpoint & Compl. \\
\midrule
Alibaba           & Qwen 2.5 7B inst    & $\mathbf{+3.06}$ & $+2.91$ & $0.99$ \\
Zhipu$^{\ast}$    & GLM 4 9B chat       & $\mathbf{+1.36}$ & $-0.10$ & $0.99$ \\
Baichuan          & Baichuan 2 7B chat  & $+0.42$          & $+0.17$ & $1.00$ \\
01.AI$^{\dagger}$ & Yi 1.5 9B chat      & $\mathbf{-0.25}$ & $-0.72$ & $1.00$ \\
\bottomrule
\end{tabular}
\caption{Provisional Chinese-side post-trained-minus-base differences and first-token compliance under the original protocol. Three of four labs shift in the positive China-favourability direction; 01.AI (Yi) shifts in the negative direction. Only Alibaba's Qwen has a clearly positive absolute endpoint under this probe. $^{\ast}$~GLM chat numbers reflect prefix-conditioned scoring (supplied prefix: newline); under na\"ive scoring A/B compliance is $3{\times}10^{-8}$. $^{\dagger}$~Yi chat numbers reflect prefix-conditioned scoring (supplied prefix: ``(''); its unmodified first-token distribution places 40\% mass on ``('' for an ``(A)''/``(B)'' response format. Because these conditions change the scoring context, full-sequence estimates from the unmodified context are required for the resubmission.}
\label{tab:chinese}
\end{table}

\subsection{Response-Format Heterogeneity at the First Token}

We observe three first-token patterns (Figure~\ref{fig:refusal}). \textbf{GLM~4 chat} places $P(\text{newline}) = 1.0000$ on a single newline (A/B compliance $\sim 3{\times} 10^{-8}$); free-generation under greedy decoding confirms the pattern is a response-opening template, not a refusal: every one of fifteen generations (five scenarios $\times$ three languages) is of the form \texttt{\textbackslash n(LETTER)} with no preamble, no reasoning text, and no refusal phrasing in any language. Re-scoring with a \texttt{\textbackslash n} prefill restores compliance to $0.91$--$0.99$ and reveals a $+1.36$ pro-China shift ($-1.47 \to -0.10$). \textbf{Yi~1.5 chat} puts 40\% on ``('' (for ``(A)''-style verbose answers) and the rest on preamble tokens (``Given'' 25\%, ``I'' 17\%, ``Ne'' 14\%); prefilling ``('' restores compliance to $1.00$ and reveals a $-0.25$ anti-China shift ($-0.47 \to -0.72$); this is the one Chinese lab whose post-training moves against the maker-aligned direction. Under na\"ive scoring the attenuated reading had reversed sign. \textbf{Qwen~2.5 inst} concentrates on ``A''/``B'' directly. GLM additionally shows a spread compression post-training ($\sigma: 19.9 \to 3.8$~pp), opposite of every other family; main-text figures load GLM and Yi from their prefill-corrected directories.

\subsection{Exploratory Free-Generation Comparison}

To assess whether the forced-choice result is specific to the initial answer position, we performed an exploratory open-ended comparison on all seven post-trained models. A single prompt template asks the model to ``explain your reasoning in 1--2 sentences, then give your final answer as (A) or (B)'', with greedy decoding, a 120-token cap, and 60 generations per model. For each generation we record the letter commit ($\pm 1$ for China) and the commit-position log-odds $\log P(\text{A}) - \log P(\text{B})$ at the answer token, conditioned on the model's own preceding text. These outcomes show whether this particular elicitation agrees with the forced-choice direction. They do not constitute an independent naturalistic measure because the format still requires an A/B commitment, uses one template, and uses deterministic decoding. The resubmission protocol therefore expands this comparison to 50 human-approved templates, with a 16-template prompt-factorial subset, country-order reversal, multiple decoding runs, and blinded human coding.

\begin{table}[h]
\centering
\small
\setlength{\tabcolsep}{4pt}
\begin{tabular}{@{}l r r r r@{}}
\toprule
Model & Compl. & Letter & Commit & Forced \\
\midrule
Mistral 7B inst    & $1.00$ & $-0.50$ & $\mathbf{-3.27}$ & $-2.77$ \\
LLaMA 3 8B inst    & $0.97$ & $\phantom{-}0.00$ & $\mathbf{-0.52}$ & $-0.96$ \\
Gemma 4 8B it      & $0.98$ & $-0.42$ & $\mathbf{-4.65}$ & $-1.01$ \\
Qwen 2.5 7B inst   & $1.00$ & $+0.03$ & $\mathbf{+0.84}$ & $+2.91$ \\
Baichuan 2 7B chat & $0.98$ & $-0.12$ & $\mathbf{-0.37}$ & $+0.17$ \\
Yi 1.5 9B chat     & $0.98$ & $-0.02$ & $\mathbf{-0.38}$ & $-0.72$ \\
GLM 4 9B chat      & $0.98$ & $-0.05$ & $\mathbf{-2.47}$ & $-0.10$ \\
\bottomrule
\end{tabular}
\caption{Exploratory single-template free-generation comparison across seven post-trained models. ``Letter'' is the mean signed letter commit ($\pm 1$ for China); ``Commit'' is the commit-position log-odds at the answer token after the model's preceding text; ``Forced'' is the no-reasoning forced-choice score for comparison.}
\label{tab:freegen}
\end{table}

All seven models commit to a parseable ``(A)''/``(B)'' at $\geq 97\%$. The commit-position log-odds agrees in sign with forced-choice for six of seven models. Baichuan is the single sign discrepancy, with both estimates within $0.4$ log-odds of the probe's zero point. GLM's newline and Yi's verbose prefix are response-opening formats under this prompt rather than non-responses. Magnitudes change materially across elicitation formats, however, so the forced-choice and open-ended measures are not interchangeable. The per-family observations below reports per-family observations without assigning them to architectural or reasoning mechanisms.

\section{Additional Methodological Detail}\label{app:methods}

\paragraph{A.1 Token-variant sums.} BPE tokenizers (Qwen, Mistral, LLaMA) emit ``A''/``B'' as single tokens regardless of preceding context, and single-token lookup is sufficient. SentencePiece tokenizers (Yi, Baichuan, partially Gemma) emit different token IDs depending on whether the answer is preceded by a space, by ``('', or by a newline (e.g. ``\_A'', ``(A'', ``\textbackslash nA''). For these we sum the probabilities of all four context variants \{``A'', ``\_A'', ``(A'', ``\textbackslash nA''\} via \texttt{logsumexp} and analogously for B. Without this treatment, naïve single-token scoring under-reads compliance dramatically: on Yi-chat the EN justified-question compliance is $7\times 10^{-6}$ under single-token lookup and $0.996$ under variant-sum, while preserving per-scenario bias at Pearson $0.9996$ between the two scorings (i.e. the variant-sum corrects the magnitude without distorting the ratio).

\paragraph{A.2 Question-polarity coherence diagnostic.} A pro-Country-X preference should produce a positive score on the \emph{justified} probe and a negative score on the \emph{unjustified} probe when X is a candidate. Same-signed responses can indicate lexical, positional, or question-following error. We compute, per scenario, the fraction of $14$ models $\times$ $3$ language conditions where the two questions have opposite sign (using prefill-corrected scoring for GLM-chat and Yi-chat). The $70\%$ threshold retains 31 scenarios; $50\%$ retains 75; and $90\%$ retains 18. Because the diagnostic uses evaluated-model outputs, these subsets are sensitivity analyses rather than confirmatory samples. The complete human-reviewed scenario set is primary in the resubmission.

\paragraph{A.3 Statistical tests.} The maker-direction analysis treats each family's post-training $\Delta$ in China-favourability as one observation ($n=7$). The two-sided exact binomial gives $k=6$ of $n=7$, $p=0.125$. A one-sample $t$-test on signed magnitude gives $t=2.78$, $p=0.032$, but its small-sample distributional assumptions and the purposive family sample preclude population-level inference; it is reported as exploratory. Within-family 95\% intervals cluster observations by scenario type. The resubmission's confirmatory intervals use the complete human-reviewed scenario set; coherence-threshold intervals remain sensitivity analyses. Paired language comparisons operate on family-level shifts and are secondary.

\paragraph{A.4 Free-generation: per-family observations.} Qwen's letter-commit average is $+0.03$ despite its forced-choice $+2.91$. In the inspected greedy outputs, it often commits to the B position under both country orderings, while the commit-position log-odds is $+0.84$ and sign-matched to forced-choice. Gemma and GLM have more negative commit-position scores ($-4.65$ and $-2.47$) than forced-choice scores ($-1.01$ and $-0.10$). Under a fixed neutral-filler condition, GLM scores $-0.32$ and Gemma $-2.35$. These context manipulations show that the score depends on text preceding the answer token. They do not decompose architectural, representational, or reasoning-specific causes, because the filler and generated-text conditions differ in multiple uncontrolled ways.

\paragraph{A.5 Pilot exclusions and scope.} Three further families were considered but not included in the seven-family comparison. \emph{(a) DeepSeek-LLM~7B} failed the pilot question-polarity validity check: per-scenario correlations between the \emph{justified} and \emph{unjustified} runs were $+0.91$ (base) and $+0.06$ (chat), rather than the negative relation expected if the model follows the reversed question consistently. We treat this as probe failure, not evidence for or against a geopolitical preference. Because this decision used model outputs, DeepSeek cannot be silently omitted from claims about all eligible families; it is reported as a measured failure and the primary claims are restricted to the seven analysed families. \emph{(b) Sub-1B models} (TinyLlama, Pythia-410M) selected a constant letter in pilot runs, so position reversal did not yield an interpretable country score. \emph{(c) InternLM~2.5~7B} could not be run because its custom modelling code was incompatible with the pinned \texttt{transformers} environment. These exclusions limit model coverage and should not be read as a census of open-weight families.

\section{Response-Format Diagnostic}\label{app:firsttoken}

\begin{figure*}[!htbp]
  \centering
  \includegraphics[width=0.95\linewidth]{figures/figure4_refusal_diagnostic.png}
  \caption{Three Chinese labs, three response patterns at the first token. Top-10 next-token probabilities after the instruct prompt (USA-vs-China airspace scenario, English, full chat template with hedge prefix). Left: GLM~4 chat places $P(\text{newline}) = 1.0000$ on a single newline, a response-opening template, not a refusal. Middle: Yi~1.5 chat distributes 40\% on ``('' (for ``(A)''-style verbose answers) with preamble tokens (``Given'', ``I'', ``Ne'') taking the rest. Right: Qwen~2.5 inst concentrates on ``A''/``B'' directly ($P(A)+P(B)=0.99$). The prefill corrections used throughout this paper restore compliance to $\geq 0.99$ for GLM (prefill \texttt{\textbackslash n}) and $1.00$ for Yi (prefill ``('').}
  \label{fig:refusal}
\end{figure*}

\section{Fictional-Name Sensitivity}\label{app:fictional}

\begin{figure*}[!htbp]
  \centering
  \includegraphics[width=0.95\linewidth]{figures/figure3_fictional_phonetic.png}
  \caption{Exploratory fictional-name sensitivity. Mean preference score per invented country name for the seven post-trained models. Four names were designed by the authors to suggest ethnolinguistic cues (Bretherland, Zhaodong, Al-Nuriyah, Korvachev); four were designed without a targeted cue (Terluna, Voskara, Drethia, Melvoni). The perceived identities were not independently validated.}
  \label{fig:fictional}
\end{figure*}
\end{document}
